\documentclass[../DoAn.tex]{subfiles}
\begin{document}

Chương này khép lại đồ án bằng ba việc: tổng kết các kết quả đạt được đối
chiếu với mục tiêu ban đầu đề ra ở Chương \ref{section:1.2}, nhìn nhận
thẳng thắn các hạn chế còn tồn tại của hệ thống ở thời điểm hiện tại, và
phác thảo các hướng phát triển tiếp theo theo ba mốc thời gian: ngắn hạn
(các việc có thể hoàn thiện trong vài tuần), trung hạn (các tính năng đòi
hỏi thiết kế lại một số thành phần), và dài hạn (các định hướng nghiên cứu
mở rộng cho các giai đoạn sau).

% ====================================================================
\section{Kết luận}
\label{section:6.1}

\subsection{Kết quả đạt được}
\label{subsection:6.1.1}

Sau toàn bộ quá trình thiết kế, hiện thực hóa và đánh giá, đề tài đã đạt
được sáu kết quả cụ thể.

\textbf{Thứ nhất, một hệ thống WAF reverse-proxy hoàn chỉnh viết bằng Go}
với rule engine 78 luật bao phủ 13 nhóm tấn công OWASP. Trên harness đánh giá
\texttt{cmd/wafeval} (46 payload đa nhóm), \textit{riêng lớp rule} chặn
\textbf{87.0\%} (40/46) và bắt được (block + challenge) \textbf{89.1\%}
(41/46); năm mẫu lọt là các trường hợp được thiết kế để lớp ML hoặc cơ chế
stateful đảm nhận (chi tiết ở mục \ref{subsection:5.2.3}). Tỷ lệ false positive
trên tập hợp lệ cố tình gây nhiễu là \textbf{3.8\%} (2/52), đáp ứng mục tiêu
\(<\) 5\%. Latency thêm vào ở phân vị 99 là 1.2 mili giây trên toàn bộ traffic,
đáp ứng mục tiêu \(<\) 5 mili giây.

\textbf{Thứ hai, tích hợp thành công lớp ML xác minh trong vùng xám}. Mô
hình DistilBERT fine-tuned 10 lớp (phiên bản v5 đang dùng production) đạt
binary attack precision \textbf{0.976} và recall \textbf{0.988} trên tập
kiểm thử out-of-distribution. Cơ chế gray-zone activation đảm bảo ML chỉ
được gọi cho 5--10\% lưu lượng, giữ overhead latency P99 toàn cục ở mức
không đáng kể.

\textbf{Thứ ba, dashboard quản trị hoàn chỉnh}. Hệ thống có 7 trang giao
diện (login, register, dashboard chính, settings cá nhân, user management,
system settings, rule builder) phục vụ ba nhóm người dùng đã được xác
định ở \ref{subsection:2.1.1}. Toàn bộ các trang đều có theme đồng nhất
với chế độ dark/light, pre-paint script chống FOUC, và glass-card layout
thống nhất.

\textbf{Thứ tư, hệ thống xác thực và phân quyền bảo mật}. JWT lưu trong
HttpOnly cookie với cờ Secure, bcrypt cost factor 10 cho mật khẩu, RBAC
ba vai trò (admin/editor/viewer) với CHECK constraint ở mức DB schema.
Endpoint \texttt{/register} công khai bịt lỗ leo thang đặc quyền (role
luôn = viewer). Các bất biến last-admin protection và self-action protection
được hiện thực hóa ở cả tầng handler lẫn tầng UI. Toàn bộ 10 test case
auth (bao gồm các regression test cho lỗ leo thang) đều pass.

\textbf{Thứ năm, pipeline huấn luyện ML có kiểm soát chất lượng}. Hard
gate bốn tiêu chí (accuracy, cmdi F1, log4shell F1, binary P/R) đã thực
sự ngăn được phiên bản v6 (accuracy 0.8866) khỏi vào production --
chứng minh giá trị thực tế của cơ chế trong việc tránh regression. Phân
tích root-cause của các trường hợp misclassification đã định hướng cụ thể
nội dung của 1800 mẫu augmentation cho v6.1.

\textbf{Thứ sáu, triển khai một lệnh qua Docker Compose}. Toàn bộ stack
(WAF + ML service + PostgreSQL) khởi động bằng một dòng lệnh
\texttt{docker compose up -d} sau khi chuẩn bị môi trường. Image Docker
được tối ưu (WAF binary tĩnh < 30 MB) phù hợp cho mọi môi trường có hỗ
trợ container.

\subsection{Đánh giá so với mục tiêu ban đầu}
\label{subsection:6.1.2}

Bảng \ref{table:goal-vs-result} đối chiếu trực tiếp từng mục tiêu đề ra ở
Chương \ref{section:1.2} với kết quả thực tế đạt được.

\begin{table}[h!]
\centering
\caption{Đối chiếu mục tiêu ban đầu với kết quả đạt được}
\label{table:goal-vs-result}
\small
\begin{tabular}{|p{5.5cm}|p{5cm}|c|}
\hline
\textbf{Mục tiêu} & \textbf{Kết quả thực tế} & \textbf{Đánh giá} \\
\hline
WAF reverse-proxy hiệu năng cao viết bằng Go &
Hoàn thành, latency P99 = 1.2 ms toàn traffic &
\checkmark Đạt \\
\hline
Bộ luật 78 rules / 13 nhóm OWASP &
78/78 rule nạp; lớp rule chặn 87.0\%, bắt 89.1\% (46 payload) &
\checkmark Đạt \\
\hline
Tích hợp ML (DistilBERT) trong vùng xám &
v5 đạt P 0.976 / R 0.988, gray-zone activation hoạt động &
\checkmark Đạt \\
\hline
Dashboard quản trị real-time &
7 trang HTML, theme đồng nhất, real-time counter &
\checkmark Đạt \\
\hline
Pipeline huấn luyện ML tái lập &
Hard gate, augmentation chiến lược, versioning rõ ràng &
\checkmark Đạt \\
\hline
Block rate \(\ge\) 95\% (mục tiêu hệ thống) &
Lớp rule: 87.0\% block / 89.1\% bắt (41/46); 5 ca biên chuyển ML &
\(\triangle\) Đạt một phần (chưa đo kèm ML) \\
\hline
False positive \(<\) 5\% &
3.8\% (2/52) trên tập đối nghịch &
\checkmark Đạt \\
\hline
Hard gate model v6 đạt accuracy \(\ge\) 0.92 &
v6 accuracy 0.8866 (fail), v5 đang dùng &
\(\triangle\) Chưa đạt với v6 \\
\hline
Rate limiting hiệu quả &
Block rate 65.6\% với burst 1500 &
\(\triangle\) Cần tune thêm \\
\hline
Triển khai qua Docker Compose &
\texttt{docker compose up -d}, image WAF \(<\) 30 MB &
\checkmark Đạt \\
\hline
\end{tabular}
\end{table}

Tổng cộng bảy trên mười mục tiêu đã được hoàn thành; ba mục tiêu còn lại
(block rate \(\ge\) 95\% ở mức hệ thống, model v6, và rate limiting tối ưu)
ở trạng thái đạt một phần với kế hoạch rõ ràng để hoàn thiện. Riêng block
rate: lớp rule đã bắt 89.1\% và năm ca biên được thiết kế để lớp ML xác minh
đảm nhận; việc đo block rate end-to-end (rule + ML) là bước đánh giá kế tiếp.

\subsection{Hạn chế}
\label{subsection:6.1.3}

Đề tài không né tránh các hạn chế còn tồn tại của hệ thống. Bốn hạn chế
chính cần được nêu rõ ràng để định hướng các bước tiếp theo và để người
sử dụng tương lai có hiểu biết chính xác về phạm vi áp dụng.

\textbf{Hạn chế 1 -- Model ML phiên bản v6 chưa đạt hard gate}. Mặc dù v5
đang phục vụ production với binary metric chấp nhận được, đường lên v6 đã
gặp một bước thoái lui không lường trước (accuracy giảm 4\% so với v5).
Hệ thống hiện chưa có một mô hình ML đạt đầy đủ năm tiêu chí gate; v6.1
đang trong giai đoạn kế hoạch với 1800 mẫu augmentation mục tiêu.

\textbf{Hạn chế 2 -- Chưa kiểm thử với traffic production thực tế}. Toàn
bộ đánh giá hiện tại dựa trên DVWA, Juice Shop, và các traffic synthetic
được sinh từ template. Hệ thống chưa được triển khai trên một ứng dụng
sản xuất với traffic thật, lưu lượng đa dạng, và nhiều phân phối payload
khác nhau. Có thể sẽ xuất hiện các vấn đề (false positive, edge case rule
không cover, vấn đề performance) chỉ lộ ra khi gặp traffic thực tế.

\textbf{Hạn chế 3 -- Rate limiting hiện chỉ đạt 65.6\% với burst 1500}.
Với một burst lớn vượt nhiều lần burst size cấu hình, vẫn có 34.4\% request
lọt qua. Mặc dù chấp nhận được với mục tiêu cản trở brute-force, kết quả
này cho thấy thuật toán Token Bucket per-IP đơn lẻ không đủ để chống các
tấn công DDoS volumetric. Cần kết hợp với lớp bảo vệ ở tầng cao hơn (CDN,
LB rate limit) hoặc nâng cấp sang thuật toán phức tạp hơn.

\textbf{Hạn chế 4 -- Phạm vi inspection chỉ tới REQUEST}. Hệ thống chỉ
kiểm tra request đầu vào, không kiểm tra response đầu ra. Điều này nghĩa
là không phát hiện được các trường hợp data leak (server vô tình trả về
SQL error chứa schema), hoặc các pattern phản ánh tấn công thành công
(stack trace, internal IP trong response body). Đây là một lớp bảo vệ
quan trọng mà các WAF thương mại đều có nhưng đề tài chưa hiện thực hóa.

% ====================================================================
\section{Hướng phát triển}
\label{section:6.2}

Dựa trên các kết quả đạt được, các hạn chế đã nêu, và quan sát về hệ sinh
thái WAF rộng hơn, đề tài đề xuất ba lộ trình phát triển theo các mốc
thời gian khác nhau.

\subsection{Hướng ngắn hạn (1--3 tháng)}
\label{subsection:6.2.1}

Các nhiệm vụ ngắn hạn tập trung vào việc hoàn thiện những gì đã bắt đầu
trong đề tài và đóng các hạn chế đã được xác định rõ.

\textbf{Hoàn thiện model v6.1}. Sinh 1800 mẫu augmentation có mục tiêu
giải quyết ba failure pattern của v6 (JSON-body bias, shell URL paths,
log4shell obfuscation -- chi tiết ở \ref{subsection:4.6.3}). Huấn luyện
v6.1 và đánh giá theo hard gate đầy đủ. Mục tiêu cụ thể: accuracy
\(\ge\) 0.92, cmdi F1 \(\ge\) 0.85, log4shell F1 \(\ge\) 0.90, binary P/R
\(\ge\) 0.95. Nếu đạt, deploy v6.1 thay thế v5 trong production.

\textbf{Cải thiện rate limiting}. Thay thuật toán Token Bucket đơn lẻ
bằng kết hợp \textit{Sliding Window Counter + Token Bucket}: Sliding
Window cho rate trung bình chính xác hơn, Token Bucket cho khả năng burst.
Đồng thời chuyển store từ in-memory sang \textbf{Redis-backed} để hỗ trợ
multi-instance deployment (các WAF instance share state qua Redis, không
bị bypass khi attacker rotate qua nhiều instance).

\textbf{Thêm RESPONSE phase inspection}. Bổ sung một bước trong middleware
stack kiểm tra response trước khi trả về client. Các pattern cần kiểm: SQL
error message (PostgreSQL/MySQL/MSSQL stack trace), file path leak
(\texttt{/etc/passwd}, \texttt{C:\textbackslash Windows}), stack trace
Java/PHP/Python với class name hệ thống. Khi phát hiện, có thể sanitize
response (thay error message generic) hoặc gửi cảnh báo cho admin.

\textbf{Tích hợp Prometheus/Grafana dashboard}. Mặc dù hệ thống đã xuất
metric ở định dạng Prometheus qua \texttt{/metrics}, chưa có Grafana
dashboard chuẩn để visualize. Bổ sung file \texttt{grafana/dashboard.json}
mô tả 8--10 panel chính: requests/s, block rate, top attack categories,
ML latency distribution, rate limiter activity, audit log volume.

\subsection{Hướng trung hạn (3--12 tháng)}
\label{subsection:6.2.2}

Các nhiệm vụ trung hạn yêu cầu thiết kế lại hoặc bổ sung các thành phần
mới với ảnh hưởng vượt ra ngoài phạm vi một tính năng đơn lẻ.

\textbf{Kiến trúc plugin: rule engine có thể mở rộng qua Lua hoặc WASM}.
Hiện tại, mọi rule đều phải tuân thủ một schema JSON cố định -- không
hỗ trợ logic tùy chỉnh phức tạp (ví dụ: ``nếu request có cookie X VÀ
header Y thì áp dụng rule Z''). Một kiến trúc plugin với engine nhúng Lua
(sử dụng \texttt{gopher-lua}) hoặc WASM (qua \texttt{wazero}) sẽ cho phép
người dùng cao cấp viết rule với điều kiện và transform tùy ý mà vẫn duy
trì isolation an toàn.

\textbf{Active learning từ false positive/negative do admin label}. Hiện
tại, audit log chỉ là dữ liệu một chiều: model đánh giá traffic, admin
xem log. Bổ sung cơ chế cho phép admin ``flag'' một bản ghi log là
\textit{false positive} (model sai khi block normal) hoặc \textit{false
negative} (model sai khi allow attack). Các bản ghi đã flag được tích lũy
thành tập huấn luyện bổ sung, định kỳ kích hoạt một lần fine-tune nhẹ.
Đây là vòng phản hồi \textit{human-in-the-loop} đặc trưng cho các hệ
thống ML production.

\textbf{Clustering mode: shared state qua Redis}. Hiện tại, mỗi WAF instance
hoạt động độc lập -- rate limiter, blacklist động, ML cache không được
share. Trong môi trường có nhiều instance phía sau load balancer, attacker
có thể rotate qua các instance để tăng burst limit hiệu quả. Chuyển các
shared state quan trọng sang Redis (rate limit bucket, IP blacklist hit
counter, ML response cache) để các instance hoạt động như một cluster
logic đồng nhất.

\textbf{Hỗ trợ HTTP/2 và HTTP/3 native}. Hiện tại, hệ thống dùng
\texttt{httputil.ReverseProxy} hỗ trợ HTTP/1.1 và HTTP/2. HTTP/3 (QUIC)
đang được triển khai rộng rãi, và một WAF không hỗ trợ HTTP/3 sẽ bị bypass
khi client và backend chuyển sang giao thức này. Thư viện
\texttt{quic-go} đã trưởng thành đủ để tích hợp.

\subsection{Hướng dài hạn (12 tháng trở lên)}
\label{subsection:6.2.3}

Các định hướng dài hạn mang tính chất nghiên cứu hơn là phát triển sản
phẩm; chúng có thể trở thành các đề tài nghiên cứu riêng biệt.

\textbf{Tích hợp threat intelligence feed}. Kết nối hệ thống với các nguồn
threat intelligence công cộng (AbuseIPDB, MISP, Spamhaus DROP list) và
thương mại (CrowdStrike, Recorded Future). IP reputation, malware URL,
known TOR exit node được tự động sync và áp dụng cho engine với policy
được cấu hình. Đây là một trong những lợi thế cốt lõi của WAF thương mại
mà giải pháp mã nguồn mở cần lấp đầy.

\textbf{Zero-day detection bằng unsupervised anomaly detection}. Bổ sung
một lớp model thứ hai song song với DistilBERT: unsupervised model
(autoencoder hoặc one-class SVM trên embedding) được huấn luyện chỉ trên
traffic normal. Bất kỳ request nào lệch xa khỏi cluster normal đều được
flag, ngay cả khi không match bất kỳ rule hay nhãn tấn công nào đã biết.
Đây là cơ chế bắt zero-day mà supervised model không thể có (vì supervised
chỉ biết các lớp đã được label).

\textbf{Compliance reporting cho PCI-DSS, GDPR, ISO 27001}. Các tổ chức
chịu sự ràng buộc của các tiêu chuẩn tuân thủ thường yêu cầu báo cáo
định kỳ về các sự kiện bảo mật, log retention theo thời gian quy định,
và evidence của các biện pháp phòng chống. Hệ thống có thể bổ sung module
sinh báo cáo theo template chuẩn của từng tiêu chuẩn, audit trail không
thể chỉnh sửa (write-once log), và policy retention tự động.

\textbf{Phân tích hành vi người dùng (User Behavior Analytics -- UBA)}.
Theo dõi pattern truy cập của từng người dùng đã đăng nhập (sequence các
endpoint, thời gian giữa request, geolocation của IP) và detect các bất
thường (đăng nhập từ vị trí mới, truy cập endpoint không thường lệ, hành
vi giống bot). UBA bổ sung một chiều bảo vệ khác cho lớp authentication
mà mật khẩu và JWT đơn lẻ không cung cấp được.

\textbf{Mở rộng sang API Gateway với schema validation}. Đối với các API
RESTful/GraphQL, hệ thống có thể đọc OpenAPI/GraphQL schema và validate
mọi request có đúng cấu trúc, kiểu dữ liệu, và phạm vi giá trị hay không.
Đây là một dạng phòng thủ chủ động vượt ra ngoài pattern matching: từ
chối request không tuân thủ schema ngay cả khi không trigger rule nào.

\bigskip
\noindent\textbf{Kết chương.}
Đồ án đã hoàn thành phần lớn các mục tiêu đề ra: một hệ thống WAF mã nguồn
mở hiệu năng cao, tích hợp đầy đủ rule engine, lớp ML xác minh, dashboard
quản trị, lớp xác thực bảo mật, và pipeline huấn luyện có kiểm soát chất
lượng. Sáu kết quả cụ thể đã được trình bày kèm số liệu định lượng. Bốn
hạn chế tồn tại được nhìn nhận thẳng thắn và đều có kế hoạch khắc phục
trong giai đoạn ngắn hạn hoặc trung hạn. Mười ba định hướng phát triển
được đề xuất, trải đều ba mốc thời gian, mở đường cho việc tiếp tục đề
tài thành một dự án nghiên cứu và phát triển dài hạn hơn. Hệ thống ở trạng
thái hiện tại đã sẵn sàng làm nền tảng cho các thử nghiệm tiếp theo trong
môi trường giáo dục và lab, và là một bước đệm hợp lý hướng tới một sản
phẩm mã nguồn mở có thể triển khai trong môi trường sản xuất sau khi hoàn
thiện các hạng mục ngắn hạn đã nêu.

\end{document}
